\section{From bounded heights to polycyclicity}
\label{sec:recognition}

We now apply Theorem~\ref{thm:bounded-height} to an arbitrary finitely
generated group quasi-isometric to a virtually polycyclic group.  The
geometric step produces a homomorphism whose kernel is locally virtually
nilpotent, meaning that every finitely generated subgroup of the kernel is
virtually nilpotent.  This proves elementary amenability.  Finiteness and
duality theorems then supply a polycyclic subgroup of finite index.

\subsection{The two Lie models}

Two Lie groups serve different purposes below.  A solvable Lie group $S$
contains a finite-index subgroup of the comparison group as a uniform
lattice; its compact quotient supplies finiteness and Poincar\'e duality.
A real-triangulable group $G$, quasi-isometric to $S$, supplies the height
geometry.  We do not require a lattice in $G$.

For the lattice model, we use the following finiteness terminology.
A group $L$ is of type $F$ if it has a finite CW classifying space.
It is of type $\FP_n(\Z)$ if the trivial module $\Z$ has a projective resolution over $\Z L$ that is
finitely generated through degree $n$, and of type $\FP_\infty(\Z)$ if
this holds in every degree.  Its integral cohomological dimension
$\operatorname{cd}_{\Z}L$ is the minimum length of a projective
$\Z L$-resolution of $\Z$, or infinity if no finite-length resolution
exists.  Type $F$ implies both $\FP_\infty(\Z)$ and finite integral
cohomological dimension.  The recognition argument will use these two
consequences separately.

We use the algebraic characterization of an integral Poincar\'e duality
group of dimension $n$: the trivial module $\Z$ admits a finite-length
resolution by finitely generated projective $\Z L$-modules, and
\[
 H^i(L;\Z L)=0\quad(i\ne n),
 \qquad H^n(L;\Z L)\cong\Z\quad\text{as abelian groups}.
\]
The right $L$-module $H^n(L;\Z L)$ is the orientation module; its action
may be nontrivial, so orientability is not part of our convention.
See \cite[Theorem~9.2 and Section~9.10]{Bieri1981} and the discussion preceding
\cite[Corollary~4.48]{Li2018}.

\begin{lemma}\label{lem:model-reduction}
Let $P$ be a finitely generated virtually polycyclic group, and let $H$ be
a finitely generated group quasi-isometric to $P$.  There exist a
finite-index subgroup $\Gamma\leq P$, a connected simply connected
solvable Lie group $S$, and a connected simply connected real-triangulable
unimodular Lie group $G$ such that
\[
 \Gamma\text{ is a uniform lattice in }S,
 \qquad H\simeq_{\mathrm{QI}}P\simeq_{\mathrm{QI}}S
             \simeq_{\mathrm{QI}}G.
\]
The group $H$ is amenable.  Moreover, $\Gamma$ has a finite classifying
space and is an integral Poincar\'e duality group of dimension $\dim S$.
\end{lemma}

\begin{proof}
Dekimpe's lattice realization theorem
\cite[Theorem~4.1]{Dekimpe2000} gives $\Gamma$ and $S$, with $\Gamma$
normal of finite index in $P$ if desired.  The uniform lattice acts
properly and cocompactly by left translations on $S$.  Thus its inclusion
is a quasi-isometry by the Milnor--\v{S}varc lemma; the same observation
applies to a finite-index subgroup acting on a Cayley graph.  See
\cite[Theorem~4.C.5(1) and Proposition~4.C.11(3)]{CornulierDeLaHarpe2016}.
Cornulier--Tessera's reduction
\cite[Lemma~3.A.1]{CornulierTessera2017} gives a real-triangulable $G$
quasi-isometric to $S$.  It remains to check unimodularity and the stated
properties of $H$ and $\Gamma$.

The group $P$ is amenable, since cyclic groups are amenable and this
property is preserved by extensions and finite extensions.  We give the
metric F\o lner argument that transfers this information to $G$ and $H$.
It is important here to obtain metric F\o lner sets in $G$, rather than
only amenability as a locally compact group; compare
\cite[Proposition~11.11 and Corollary~11.13]{Tessera2008}.

Let $Y$ denote either $G$, with a left invariant Riemannian metric and
left Haar measure $\mu$, or $H$, with a word metric and counting measure.
Choose a quasi-isometry $u:P\to Y$ with coarse-surjectivity radius $c>0$,
and fix $a>c$.  For finite nonempty $E\subset P$, put
\[
 A_E=\bigcup_{p\in E} B_Y(u(p),a).
\]
For each fixed radius, the lower quasi-isometry inequality bounds the
number of points of $P$ whose images lie in a ball of that radius.
Consequently the balls in this union have bounded multiplicity, and
\begin{equation}\label{eq:recognition-folner-lower}
 \mu(A_E)\geq c_1|E|
\end{equation}
for a constant $c_1>0$ independent of $E$.

Write $N_R(A)$ for the metric $R$-neighborhood of a set $A$.  If
$y\in N_R(A_E)\setminus A_E$, choose $p_y\in P$ with
$d_Y(y,u(p_y))\leq c$ and $p\in E$ with
$d_Y(y,u(p))\leq R+a+1$.  Since $a>c$, we have $p_y\notin E$.
The lower quasi-isometry inequality gives
$d_P(p_y,E)\leq R'$, with $R'$ depending only on $R,a$ and the
quasi-isometry constants.  It follows that
\begin{equation}\label{eq:recognition-folner-boundary}
 \mu\bigl(N_R(A_E)\setminus A_E\bigr)
 \leq \mu(B_Y(1,c))\,|N_{R'}(E)\setminus E|.
\end{equation}
An increasing-radius F\o lner sequence in $P$, together with
\eqref{eq:recognition-folner-lower}--\eqref{eq:recognition-folner-boundary},
therefore gives finite-measure sets $A_j\subset Y$ satisfying
\begin{equation}\label{eq:recognition-metric-folner}
 \frac{\mu(N_R(A_j))}{\mu(A_j)}\longrightarrow1
 \qquad\text{for every fixed }R.
\end{equation}
For the discrete group $H$, this is the F\o lner criterion for
amenability.

For $Y=G$, right translation by $g$ moves every point a distance
$d_G(1,g)$, so $A_jg\subset N_{d_G(1,g)}(A_j)$.  If right translation
multiplied left Haar measure by a factor greater than one, this inclusion
would contradict \eqref{eq:recognition-metric-folner}.  A nontrivial
modular homomorphism has such a value, after replacing $g$ by $g^{-1}$
if necessary.  Thus $G$ is unimodular.

Finally, a connected simply connected solvable Lie group of dimension
$n$ is diffeomorphic to $\R^n$; explicit coordinates are given in
\cite[Lemma~3.B.2]{CornulierTessera2017}.
The compact smooth manifold $\Gamma\backslash S$ is therefore
aspherical and has a finite CW structure.  It is a finite classifying
space for $\Gamma$, and manifold Poincar\'e duality makes $\Gamma$ an
integral Poincar\'e duality group of dimension $n$.
\end{proof}

\subsection{Polynomial growth in a height kernel}

The polynomial packing estimate needed here is stronger than an estimate
of the form $e^{o(t)}$.  Lemma~\ref{lem:polynomial-path-packing} gives
that estimate for paths whose heights remain bounded at every point.
We record carefully how it applies to a subgroup with its own word
metric.

\begin{lemma}\label{lem:height-kernel-growth}
Let $G$ be a connected simply connected real-triangulable unimodular Lie group with the height
homomorphism $\pi:G\to W$, and let $f:H\to G$ be a quasi-isometry from
a finitely generated group, normalized by $f(e)=1$.  Suppose $H_0\leq H$
and $\tau:H_0\to W$ is a homomorphism such that, for some $D<\infty$,
\begin{equation}\label{eq:recognition-corrected-height}
 \|\pi f(h)-\tau(h)\|\leq D\qquad(h\in H_0).
\end{equation}
Then every finitely generated subgroup of $K_0=\ker\tau$ has polynomial
growth and is virtually nilpotent.
\end{lemma}

\begin{proof}
Fix a finitely generated subgroup $J\leq K_0$ and a finite symmetric
generating set $T$ of $J$.  The heights of all orbit points $f(j)$,
$j\in J$, have norm at most $D$.  For a word
$j=t_1\cdots t_m$ with $m\leq t$, put
$j_i=t_1\cdots t_i$.  The upper quasi-isometry inequality gives a
constant $L_T$ such that
\[
 d_G(f(j_{i-1}),f(j_i))\leq L_T
 \qquad(1\leq i\leq m).
\]
Join these successive points by minimizing geodesics.  The resulting
path starts at $1$, ends at $f(j)$, and has length at most $L_Tt$.
The homomorphism $\pi$ is Lipschitz for a left invariant Riemannian
metric: its differential has constant norm in left-trivialized tangent
spaces.  If $L_\pi$ is a Lipschitz constant, every point of the path has
height norm at most $D+L_\pi L_T$.

Lemma~\ref{lem:polynomial-path-packing}, with these fixed length and
height bounds, shows that a maximal $1$-separated subset of
$f(B_T(t))$ has at most $C_T(1+t)^{b_T}$ points.  The unit balls about
that subset cover $f(B_T(t))$.  The lower quasi-isometry inequality
bounds the number of elements of $H$ mapping into each such ball by
one constant: any two such elements are at uniformly bounded distance
in the fixed word metric of $H$.  Hence
\[
 |B_T(t)|\leq C'_T(1+t)^{b_T}.
\]
This is an intrinsic growth bound for $J$; no undistortion of $J$ in $H$
is assumed.  Gromov's polynomial-growth theorem
\cite[Main Theorem, p.~54]{Gromov1981} makes $J$ virtually nilpotent.
\end{proof}

The packing input works with lifts to $N\rtimes D$:
compactly bounded increment lifts place chosen endpoint lifts in a
polynomial-volume coordinate region.  Separation in $G$ separates
these lifts, as in Lemma~\ref{lem:polynomial-path-packing}, so the
argument applies even when $N\cap D\ne\{1\}$.

\subsection{The transferred action and elementary amenability}

We use the term \emph{elementary amenable} for the smallest class of
groups containing finite and abelian groups and closed under subgroups,
quotients, extensions, and directed unions.  In the next proposition
amenability is a hypothesis, but elementary amenability is a conclusion.

\begin{proposition}\label{prop:recognition-from-height}
Let $G$ be a connected simply connected real-triangulable unimodular Lie
group with height homomorphism $\pi:G\to W$.  Let $H$ be a finitely
generated amenable group quasi-isometric to $G$.  Assume the uniform
bounded-height conclusion of Theorem~\ref{thm:bounded-height} for the
family of self-quasi-isometries transferred from left translations of
$H$.  Then there is a finite-index normal subgroup $H_0\leq H$ and a
homomorphism $\tau:H_0\to W$ whose kernel is locally virtually
nilpotent.  In particular, $H$ is elementary amenable.
\end{proposition}

\begin{proof}
Choose quasi-inverse maps $f:H\to G$ and $r:G\to H$.  Bounded changes
allow us to arrange $f(e)=1$ and $r(1)=e$.  For $h\in H$, define
\[
 F_h=f\circ L_h\circ r,
 \qquad L_h(k)=hk.
\]
The maps $F_h$ have uniform quasi-isometry constants, because the maps
$L_h$ are isometries.  There is a constant $D_0$ such that
\begin{equation}\label{eq:recognition-quasiaction}
 \sup_{g\in G}d_G(F_hF_k(g),F_{hk}(g))\leq D_0
 \qquad(h,k\in H),
\end{equation}
and $F_e$ is uniformly close to the identity.  These assertions follow
from the bounded errors in $r\circ f$ and $f\circ r$.  Also
$F_h(1)=f(h)$.

Put $s(h)=\pi f(h)$.  The assumed height theorem gives a single
$B<\infty$ and matrices $A_h$ in the fixed finite group $\A_G$
such that
\begin{equation}\label{eq:recognition-affine-action}
 \|\pi F_h(g)-s(h)-A_h\pi(g)\|\leq B
 \qquad(h\in H,\ g\in G).
\end{equation}
Apply this formula twice to $F_hF_k(g)$, and compare with its application
to $F_{hk}(g)$ using \eqref{eq:recognition-quasiaction}.  It follows
that
\[
 (A_hA_k-A_{hk})\pi(g)
\]
is bounded as $g$ ranges over $G$.  Since $\pi$ is surjective onto the
vector space $W$, its linear coefficient vanishes.  Similarly $A_e=I$.
Thus $h\mapsto A_h$ is a homomorphism, and
\[
 H_0=\{h\in H:A_h=I\}
\]
is normal and has finite index.  Evaluating the same comparison at
$g=1$ gives a constant $D_1$ such that
\begin{equation}\label{eq:recognition-defect}
 \|s(hk)-s(h)-s(k)\|\leq D_1\qquad(h,k\in H_0).
\end{equation}
The uniform constant in the height theorem is essential here: $D_1$
does not depend on $h$ or $k$.

The subgroup $H_0$ is amenable.  Let $m$ be a left-invariant mean on
bounded real functions on $H_0$, extended coordinatewise to functions
with values in $W$.  Define
\begin{equation}\label{eq:recognition-mean}
 \tau(h)=m_x\bigl(s(hx)-s(x)\bigr),\qquad h\in H_0.
\end{equation}
For fixed $h$, the function being averaged is bounded, since
\eqref{eq:recognition-defect} bounds its distance from $s(h)$ by $D_1$.
Positivity of the mean gives
\begin{equation}\label{eq:recognition-mean-error}
 \|\tau(h)-s(h)\|\leq D_1.
\end{equation}
To verify additivity, write
\[
 s(hkx)-s(x)=s(hkx)-s(kx)+s(kx)-s(x).
\]
Left invariance under $x\mapsto kx$ shows that the means of the two
terms on the right are $\tau(h)$ and $\tau(k)$.  Therefore
$\tau(hk)=\tau(h)+\tau(k)$.

Let $K_0=\ker\tau$.  Equation~\eqref{eq:recognition-mean-error} is
the hypothesis of Lemma~\ref{lem:height-kernel-growth}; that lemma
makes each finitely generated subgroup of $K_0$ virtually nilpotent.
Nilpotent groups are elementary amenable by their central series, and
so are their finite extensions.  The finitely generated subgroups of
$K_0$ form a directed family with union $K_0$.  Thus $K_0$ is elementary
amenable.  The exact sequence
\[
 1\longrightarrow K_0\longrightarrow H_0
   \overset{\tau}{\longrightarrow}\tau(H_0)\longrightarrow1
\]
and abelianness of $\tau(H_0)$ show that $H_0$ is elementary amenable.
Finally $H/H_0$ is finite, so $H$ is elementary amenable as well.
\end{proof}

The image $\tau(H_0)$ need not be discrete in $W$, and the kernel $K_0$
need not be finitely generated.  Neither property was used: abelianness
of the image and the directed-union argument are sufficient.

\subsection{Finiteness, duality, and the conclusion}

The previous proposition supplies elementary amenability.  To complete
recognition, we first transfer $\FP_\infty(\Z)$ from the comparison
lattice.  The elementary amenable finiteness theorem then gives a
finite-index subgroup of type $F$, which supplies the finite
cohomological dimension needed for duality transfer.

\begin{proof}[Proof of Theorem~\ref{thm:recognition}]
Let $P$ and $H$ satisfy the hypotheses of the theorem.  Apply
Lemma~\ref{lem:model-reduction} and retain both the real-triangulable
unimodular model $G$ and the uniform lattice $\Gamma\leq S$.

If $W=0$, Lemma~\ref{lem:canonical-height} makes $G$ nilpotent.  Its metric
balls have polynomial Haar volume, by the nilpotent coordinate bounds
of Lemma~\ref{lem:d-scaling}.  A quasi-isometry from $H$ to $G$ has
bounded multiplicity on bounded scales, so disjoint small-ball packing
transfers this to polynomial word growth of $H$.  Gromov's theorem
\cite[Main Theorem, p.~54]{Gromov1981} makes $H$ virtually nilpotent,
hence virtually polycyclic.  This also includes the zero-dimensional
model: in that case $H$ is finite.  We may therefore suppose $W\ne0$.

By Lemma~\ref{lem:model-reduction}, $H$ is amenable.  Apply
Theorem~\ref{thm:bounded-height} to the uniform family transferred
from its left translations.  Proposition~\ref{prop:recognition-from-height}
shows that $H$ is elementary amenable.

The group $\Gamma$ has a finite classifying space and is quasi-isometric
to $H$.  Li's coarse invariance theorem
\cite[Corollary~4.47]{Li2018}, applied over $\Z$ for each $n$, gives
$H\in\FP_\infty(\Z)$.  The elementary amenable finiteness theorem
\cite[Theorem~1.1, (viii)$\Rightarrow$(v)]{KrophollerMartinezPerezNucinkis2009}
now gives a finite-index subgroup $V\leq H$ of type $F$.  The structural
conclusion recorded in Section~2 of the same paper also gives a
finite-index solvable subgroup $U\leq H$.  Set $H_1=U\cap V$.
Then $H_1$ is solvable and has finite index in $H$.  The finite cover of
a finite classifying space for $V$ corresponding to $H_1$ is a finite
classifying space for $H_1$.  In particular,
\[
 \operatorname{cd}_{\Z}H_1<\infty.
\]
This also implies that $H_1$ is torsion-free: restriction of a finite
projective resolution to a hypothetical nontrivial finite cyclic
subgroup would give finite integral cohomological dimension for that
subgroup, which is impossible.

Both $H_1$ and $\Gamma$ now have finite integral cohomological dimension,
and they are quasi-isometric.  Since $\Gamma$ is an integral
Poincar\'e duality group, Li's duality transfer
\cite[Corollary~4.48]{Li2018}, with coefficient ring $\Z$, makes $H_1$
an integral Poincar\'e duality group.  Here the type-$F$ subgroup supplies
the finite-cohomological-dimension hypothesis.  Bieri's theorem
\cite[Theorem~9.23]{Bieri1981} says that a solvable integral
Poincar\'e duality group is polycyclic.  Hence $H_1$ is polycyclic, and
$H$ is virtually polycyclic.
\end{proof}

\begin{proof}[Proof of Corollary~\ref{cor:lattice-recognition}]
Let $\Lambda$ be a lattice in a connected simply connected solvable Lie
group $S$.  Such a lattice is polycyclic
\cite[Proposition~3.4]{FarbWeinberger2008}, and is uniform
\cite[Theorem~2.9(i)]{Geng2015}.  It is therefore finitely generated and
quasi-isometric to $S$ by the Milnor--\v{S}varc lemma.  A finitely
generated group $H$ quasi-isometric to $\Lambda$, or equivalently to
$S$, is virtually polycyclic by Theorem~\ref{thm:recognition}.
Dekimpe's realization theorem \cite[Theorem~4.1]{Dekimpe2000} then
embeds a finite-index subgroup of $H$ as a uniform lattice in some
connected simply connected solvable Lie group.

Conversely, every virtually polycyclic comparison group has such a
finite-index lattice model by the same realization theorem.  Finite
index preserves quasi-isometry, so the lattice formulation implies the
group formulation.  The ambient Lie group obtained for $H$ is not
prescribed by this equivalence.
\end{proof}
